\documentclass[11pt,a4paper]{article}
\usepackage[margin=1in]{geometry}
\usepackage{amsmath,booktabs,hyperref,float}

\title{Iteration 012: Closed-Form + FIR Ensemble}
\author{Autoresearch Loop}
\date{April 7, 2026}

\begin{document}
\maketitle

\begin{abstract}
A weighted average of closed-form and FIR predictions achieves $r=0.375$,
not improving over the FIR alone ($r=0.376$). The two models make highly
correlated predictions, so ensembling provides minimal complementary information.
\end{abstract}

\section{Hypothesis}
The closed-form solution is globally optimal for MSE but has no temporal context.
The FIR model adds temporal dynamics but may overfit. An ensemble could combine
the robustness of CF with the expressiveness of FIR.

\section{Method}
Weighted average: $\hat{y} = \alpha \cdot \hat{y}_\text{CF} + (1-\alpha) \cdot \hat{y}_\text{FIR}$,
with $\alpha$ optimized on the validation set via grid search.

\section{Results}
\begin{table}[H]
\centering
\begin{tabular}{lrrr}
\toprule
Model & Mean $r$ & Std $r$ & SNR (dB) \\
\midrule
FIR + dropout (iter011) & 0.376 & 0.076 & 0.61 \\
CF + FIR ensemble & 0.375 & 0.076 & 0.61 \\
\bottomrule
\end{tabular}
\end{table}

\section{Analysis}
Since the FIR model is initialized from the CF solution and remains close to it,
the two models produce highly correlated outputs. Ensembling is only beneficial
when constituent models have diverse error patterns. This suggests the FIR model's
improvement over CF is primarily from temporal refinement, not a fundamentally
different spatial mapping.

\section{Next}
Iteration 013: band-specific models to exploit frequency-dependent spatial patterns.

\end{document}
